\documentclass[11pt]{article}
\usepackage[a4paper,margin=25mm]{geometry}
\usepackage{amsmath,amssymb,amsthm,mathtools,hyperref,booktabs}
\hypersetup{colorlinks=true,linkcolor=blue,urlcolor=blue,citecolor=blue}
\newcommand{\MA}{\operatorname{MA}}
\newcommand{\NA}{\mathrm{NA}}
\newcommand{\PSH}{\operatorname{PSH}}
\newcommand{\Ent}{\operatorname{Ent}}
\newcommand{\Ric}{\operatorname{Ric}}
\newcommand{\Vol}{\operatorname{Vol}}
\newcommand{\ord}{\operatorname{ord}}
\newcommand{\cE}{\mathcal E}
\newcommand{\ddc}{dd^c}
\newtheorem{theorem}{Theorem}
\newtheorem{lemma}[theorem]{Lemma}
\newtheorem{proposition}[theorem]{Proposition}
\newtheorem{corollary}[theorem]{Corollary}
\theoremstyle{remark}\newtheorem{remark}[theorem]{Remark}
\title{Model recovery, Legendre mass profiles, and endpoint defects\\
Partial results for the original P19.2 approximation problem}
\author{}\date{6 October 2026}
\begin{document}\maketitle
\begin{abstract}
The original residual P19.2 target remains unresolved. Starting from
the precise model energy--entropy recovery theorem supplied with the
problem, we prove quantitative convergence of actual maximal rays,
integrated Legendre mass profiles and every interior level mass. An
ordered-envelope criterion gives convergence of interior singularity types
in the probability-normalized singularity pseudometric. We prove a
joint approximation/large-time entropy limit for these model rays
from a fixed smooth K\"ahler initial potential. A relatively ample
blowup of $\mathbf P^1\times\mathbf P^1$ provides a positive
discrepancy calibration and an explicit endpoint defect for actual
model sequences. Finally an energy-defect estimate explains exactly
what the supplied polar Cantor example obstructs. None of these
statements proves simultaneous preservation of every prescribed
$P$-envelope, endpoint and initial singularity signature in the
unchanged original residual target. No novelty or priority is claimed.
\end{abstract}

\section{Problem, normalization, and scope}
The original P19.2 asks for model approximations preserving the required
energy, entropy and finer prescribed singularity data in the residual
hybrid programme. It includes the singularity/envelope levels attached
to log-discrepancy-model rays. Published projective Fujita regimes are
established subcases. A counterexample to an arbitrary approximation
sequence, or a theorem for a stronger restricted input signature, is
not a resolution of this existence problem. The catalogue does not
choose one topology or an explicit list of levels for all residual
inputs. We consequently distinguish exact preservation, convergence
of level envelopes, convergence in a relative singularity topology,
and convergence of level masses. We do not identify those requirements.

Let $X$ be a connected smooth compact K\"ahler manifold of dimension
$n\geq1$, with K\"ahler form $\omega$, $\alpha=[\omega]$,
$V=\int_X\omega^n$, and reference probability $\mu_0=\omega^n/V$.
The analytic non-Archimedean carrier is the compact complex-space
valuative space $X^\beth$ over $\mathbf C$ with its trivial valuation;
semivaluation supports are retained. Throughout,
\begin{equation}\label{eq:conventions}
 \ddc=\frac{i}{2\pi}\partial\bar\partial,
 \qquad t=-\log|\tau|^2,
 \qquad \ddc\log|f|^2=[\operatorname{div}(f)].
\end{equation}
Write $\omega_u=\omega+\ddc u$ and use
\[
 \MA(u)=V^{-1}\langle\omega_u^n\rangle,
 \quad E(u)=\frac1{(n+1)V}\sum_{k=0}^n
  \int_Xu\,\omega_u^k\wedge\omega^{n-k},\quad E(0)=0.
\]
On $\cE^1(\omega)$ the measure has probability mass one.
We use the entire finite-energy extension of $E$ and the probability
$d_1$ metric; in particular $E(u+c)=E(u)+c$ and, if $u\geq v$,
$d_1(u,v)=E(u)-E(v)$.
Entropy is the full probability entropy
\[
 H(u)=\Ent(\MA(u)\mid\mu_0)=\int f\log f\,d\mu_0
\]
if $\MA(u)=f\mu_0$, and $+\infty$ otherwise. Set
\[
 H_\alpha(\phi)=\int_{X^\beth}A_X\,\MA_\alpha(\phi),
 \qquad A_X(v_{\rm triv})=0,
\]
with the probability non-Archimedean operator.

An admissible model in this paper is a smooth compact equivariant
cohomological configuration $\rho:\mathcal Y\to X\times\mathbf P^1$
and a real vertical divisor $D$ such that
$\rho_X^*\alpha+D$ is relatively K\"ahler. Its usual datum is
$f_D(v_F)=\ord_F(D)/b_F$, $b_F=\ord_F(\tau)$; the compatible
total potential satisfies $U+\log|s_D|^2=O(1)$.
For $\phi\in\cE^{1,\NA}(\alpha)$ and $w\in\cE^1(\omega)$,
$r(w,\phi)$ denotes the actual maximal weak ray, maximal with respect
to this usual datum and this actual initial point.
This is different from $P$-maximality of a Legendre test curve.

For a ray $r$ from zero put
\[
 \chi_r(\lambda)=\inf_{t\geq0}(r_t-t\lambda),\qquad
 M_r(\lambda)=V^{-1}\int_X\langle\omega_{\chi_r(\lambda)}^n\rangle.
\]
At a level identically $-\infty$ we assign the raw mass $M_r=0$.
All other masses lie in $[0,1]$. They are not probability measures
unless further normalized, and zero mass is never divided out.
$P_\omega[u]=(\lim_{C\to\infty}P_\omega(u+C,0))^*$.
The test curves of actual weak rays from zero are $P$-model:
$P_\omega[\chi_r(\lambda)]=\chi_r(\lambda)$.
For a smooth K\"ahler initial potential $w$, use background
$\omega_w$ and ray $r_t-w$. This is an actual ray from zero in the
same class. Its levels correspond, after adding $w$, to the rooftop
envelopes relative to $w$ in the original background. No such
replacement by a smooth background is asserted for a singular $w$.

\section{Exact mathematical inputs and their applications}\label{sec:inputs}
The following inputs, rather than a historical review verdict, are used.
They are stated in the signature needed here. Primary text units and
remaining source boundaries are recorded in the application ledger.

\paragraph{R: supplied model recovery.}
For every smooth compact K\"ahler $(X,\alpha)$ and
$\phi\in\cE^{1,\NA}(\alpha)$ with $H_\alpha(\phi)<\infty$, there
exist the actual smooth relatively K\"ahler configurations just
defined, with model data $\eta_j$, such that
\begin{equation}\label{eq:R}
 \eta_j\longrightarrow\phi\text{ strongly},\qquad
 H_\alpha(\eta_j)\longrightarrow H_\alpha(\phi).
\end{equation}
This is the new usable mathematical result supplied in the frozen
input, Section 5 of the supplied P19.1 paper. It concerns the usual
datum, a fixed smooth manifold and a K\"ahler class, not arbitrary
extra $P$-singularities or a degenerating family. Its proof passes
through PL obstacles, vertical special Fujita data and a locally
proved analytic Rees discrepancy bound. We use this supplied theorem
as a mathematical premise; its full construction is not claimed as a
new independent proof here.

For clarity about its actual models, its obstacle step has
$\Gamma=\rho_X^*\alpha+D$ big on an $(n+1)$-fold and uses smooth
equivariant modifications $p_k$, effective vertical $E_k$, and
$\beta_k=p_k^*\Gamma-E_k$ K\"ahler with
\[
 \beta_k^{n+1}\to\Vol(\Gamma),\quad
 \beta_k^nK_{\mathcal Y_k/\mathcal X}\to0,\quad
 \beta_k^np_k^*\xi\to\langle\Gamma^n\rangle\xi
\]
for each fixed smooth closed real $(1,1)$ class $\xi$.
Its resulting model is $f_{p_k^*D-E_k}$, followed by an entropy/strong
diagonal over obstacles. The entropy gate is the canonical-error
limit, using $0\leq\Delta_k\leq K_{\mathcal Y_k/\mathcal X}$;
volume convergence alone is insufficient. The fresh primary
Trusiani text \cite{T} is the transcendental paper of 25 September
2026, not the catalogue's polarized paper. The detailed vertical
construction and local discrepancy repair are part of the supplied
premise R; we do not promote a source abstract into their proof.

\paragraph{F: finite-energy foundations.}
Mesquita--Piccione--Witt Nystr\"om \cite{MPWN}, Lemma 5.0.5,
Theorem 5.2.2 and Theorem B, imply, in this same smooth compact
K\"ahler signature, that strong $\phi_j\to\phi$ gives convergence of
$E_\alpha$, of the suprema and of the Monge--Amp\`ere measures, and
\[
 \int|\phi_j-\phi|\,\MA_\alpha(\psi)\to0
 \quad(\psi\in\cE^{1,\NA}(\alpha)\text{ fixed}).
\]
The supporting inequality is
$E_\alpha(a)-E_\alpha(b)\leq\int(a-b)\MA_\alpha(b)$.
Analytically the corresponding supporting inequality, affine energy,
ordered $d_1$ identity, and $I_1(a,b)\leq C d_1(a,b)$ hold on
$\cE^1(\omega)$ \cite{D}. These apply to maxima of finite-energy
potentials as well. Mesquita--Piccione \cite{MP}, Theorem 5.1.7,
gives actual maximal rays, order in both arguments, and calibrated
energy $E(r(w,\phi)_t)=E(w)+tE_\alpha(\phi)$.
Lemma 5.3.1 gives $\sup_Xr(0,\phi)_t=t\sup\phi$.
The same source gives nonnegative lower semicontinuous discrepancy
$A_X$ as the directed supremum of the continuous PL retraction tests
$A_X\circ r_{\mathcal Y}$.

\paragraph{L: Legendre foundations.}
In an arbitrary big class on a smooth compact K\"ahler manifold,
Darvas--Xia--Zhang \cite{DXZ}, Theorem 2.6 and formula (51) of the
published text (formula (62) in arXiv v3), give the finite-energy
$P$-maximal-test-curve/weak-ray correspondence and
\begin{equation}\label{eq:legendre-energy}
 c_r=\lambda_r^++\int_{-\infty}^{\lambda_r^+}(M_r(\lambda)-1)d\lambda,
 \qquad c_r=E(r_1)-E(r_0).
\end{equation}
Here the initial point is the minimal potential, which in our
K\"ahler background is zero. Their relative theorem for other
positive-mass starting points has a separate signature; it is not
silently substituted in this formula. For the K\"ahler case the proof
of the energy formula is the dyadic maximum/energy-sandwich proof of
Darvas--Xia \cite{DX}, Theorem 3.7(iv), Proposition 3.8 and Lemma 3.9:
bounded test curves are handled by Riemann sums, and finite-energy
curves by decreasing bounded-ray approximation. The complete
nonpolarized proof is also \cite{DZ}, Theorems 3.7 and 3.9,
Propositions 3.8 and 3.10, and Lemma 3.11 (whose K\"ahler proof
is the cited \cite{DX}, Lemma 3.9). Only these analytic
proof units are used, not a polarization-dependent algebraic closure
theorem. Level masses are positive and continuous on
$(-\infty,\lambda_r^+)$ and nonincreasing everywhere.
For the initial-zero maximal rays, $\lambda_r^+=\sup\phi$.

For ordered singularity types, the singularity pseudometric, with
probability-normalized energy, obeys
\begin{equation}\label{eq:ordered-dS}
 d_{\mathcal S}([u],[v])=\frac1{(n+1)V}\sum_{k=0}^n
 \left(\int\omega^k\wedge\langle\omega_v^{n-k}\rangle
 -\int\omega^k\wedge\langle\omega_u^{n-k}\rangle\right)
 \quad\text{if }[u]\leq[v].
\end{equation}
This is \cite{DDNLmetric}, Lemma 3.4 and its energy formula; its
triangle inequality comes from its definition by chordal rays.
At positive mass its zero classes coincide with equal $P$-envelopes.
We also use \cite{DDNLmetric}, Lemma 4.3: for $u\leq v$ with
positive raw mass $m(u)$, and $b>1$ such that
$b^n m(u)>(b^n-1)m(v)$, the envelope
$P_\omega(bu+(1-b)v)$ is an $\omega$-psh function.
The source proves this by truncating $u$ against $v$, a contact-set
measure estimate for that envelope, and a mass contradiction;
the strict inequality and positive-mass hypotheses are retained.
In their $i\partial\bar\partial$ convention put the potential equal
to $u/(2\pi)$ and multiply the energy metric by $2\pi/V$ to get
\eqref{eq:ordered-dS}. The positive-mass envelope theorem
\cite{DDNL}, Section 3.2 and Theorem 3.12, says that for $u\leq0$
with positive non-pluripolar mass, $P_\omega[u]\geq u$ preserves
that mass and is idempotent. Zero-mass levels are not covered by
this particular theorem.

\paragraph{Q: calibrated mixed estimates and model regularity.}
We use the following exact consequence of \cite{MP}, Lemma 5.3.2
and the proof of Theorem 5.3.4. For mean-normalized finite-energy
slots, a mixed-pairing difference is bounded by
$C\varepsilon^q R^{1-q}$, $q=2^{-n}$, where $\varepsilon$ is the
maximum slot energy distance and $R$ bounds the slots' $J$ plus the
fixed supremum-minus-mean bound. The constant depends only on the
fixed smooth forms and dimension. Signed forms and PL weights are
split into positive inputs on one fixed common model.
Slots are replaced one at a time; this is not a bound depending on
the changing model of $\eta_j$.
For a fixed PL test $f_C$, a smoothly divisor-compatible weight
$\Psi_t$ has
\begin{equation}\label{eq:pairing-slope}
 \lim_{t\to\infty}t^{-1}\int\Psi_t\MA(r(0,\phi)_t)
       =\int f_C\MA_\alpha(\phi).
\end{equation}
Fixed smooth contracted energies also have their calibrated
non-Archimedean slopes, strongly continuous in $\phi$.
The calibration is proved in Section \ref{sec:mixed} below.
Sj\"ostr\"om Dyrefelt \cite{SD}, Lemma 4.6 and Remark 4.7,
give the actual model ray from any fixed smooth K\"ahler potential;
its divisor-removed total potential has locally bounded complex
Hessian. We use this regularity, not a printed entropy slope constant.

\paragraph{M: Mabuchi convexity.}
The full probability entropy with our operator has the expression
\begin{equation}\label{eq:M}
 M=H+G,\quad G=T_{-\Ric(\omega)}+\bar S E,\quad
 T_\xi(u)=\frac1V\sum_{k=0}^{n-1}
 \int u\,\xi\wedge\omega_u^k\wedge\omega^{n-1-k},
\end{equation}
where $\Ric(\omega)=-\ddc\log\omega^n$ and
$\bar S=nV^{-1}\int\Ric(\omega)\wedge\omega^{n-1}$.
This $M$ is $d_1$ lower semicontinuous and convex on canonical
finite-energy weak geodesics; smooth contracted energies are $d_1$
continuous. This is \cite{BDL}, Theorem 4.7 with zero twist and
Propositions 4.4--4.5. To apply its $i\partial\bar\partial$ version,
use $u/(2\pi)$ with the same form $\omega$. Currents and probability
measures coincide. The Ricci and average scalar curvature scale by
$2\pi$, cancelling the energy scaling; the entropy coefficient is
one. Differentiating $H$ directly gives
$nV^{-1}\int\dot u(\Ric\omega-\Ric\omega_u)\wedge\omega_u^{n-1}$,
which checks \eqref{eq:M}.

All these applications are on a fixed manifold. No compactness
theorem for genuinely degenerating families or equivariant stability
correspondence is imported.

The bounded local products used below are the Bedford--Taylor
products: locally bounded psh slots have symmetric multilinear
products, continuous under decreasing locally bounded limits
\cite{BT}, Theorem 2.1 and Corollary 2.2. For locally $C^{1,1}$
potentials their products are the almost-everywhere coefficient
products. Indeed local convolution has bounded Hessians, converging
at almost every Lebesgue point and in every finite $L^p$; the
polynomial coefficient products converge in $L^1$, and the
Bedford--Taylor local continuity identifies this limit with the
current product. Smooth local backgrounds may be added throughout.

\section{A common upper ray and the mass-profile estimate}
\begin{theorem}\label{thm:profiles}
Let $\eta_j\to\phi$ strongly in $\cE^{1,\NA}(\alpha)$. Models
and finite entropy are not needed for this theorem. Put
$r_j=r(0,\eta_j)$, $r=r(0,\phi)$,
$\lambda_j=\max(\eta_j,\phi)$, and
\[
 B_j=2E_\alpha(\lambda_j)-E_\alpha(\eta_j)-E_\alpha(\phi).
\]
Then $B_j\geq0$, $B_j\to0$, and for every $t\geq0$,
\begin{equation}\label{eq:ray-bound}
 d_1(r_{j,t},r_t)\leq tB_j.
\end{equation}
The chordal distance tends to zero and
\begin{equation}\label{eq:mass-bound}
 \int_{\mathbf R}|M_{r_j}-M_r|\,d\lambda\leq B_j.
\end{equation}
In particular $M_{r_j}(\lambda)\to M_r(\lambda)$ at every
$\lambda<\sup\phi$, uniformly on each compact subinterval of
$(-\infty,\sup\phi)$. These are raw masses of the actual
$P$-model levels. No equality of envelopes at each index is asserted.
\end{theorem}
\begin{proof}
The maximum $\lambda_j$ is psh and finite energy. The supporting
inequality and F give
\[
 0\leq E_\alpha(\lambda_j)-E_\alpha(\phi)
 \leq\int|\eta_j-\phi|\MA_\alpha(\phi)\to0.
\]
Together with $E_\alpha(\eta_j)\to E_\alpha(\phi)$ this proves the
assertions about $B_j$. The actual ray $s_j=r(0,\lambda_j)$ dominates
both rays. Thus triangle inequality and ordered energy give
\[
 d_1(r_{j,t},r_t)\leq d_1(r_{j,t},s_{j,t})+d_1(s_{j,t},r_t)
 =tB_j.
\]
It is $s_j$, not the pointwise maximum of two rays, that is used.

Legendre transforms preserve order. Hence the level curves
$\sigma_j=\chi_{s_j}$ dominate $\chi_{r_j}$ and $\chi_r$.
Non-pluripolar mass monotonicity implies
$M_{s_j}\geq M_{r_j},M_r$. All three test curves have finite energy.
Choose any finite $L$ above their three cutoffs. Formula
\eqref{eq:legendre-energy} can be written
\[
 c_a=L+\int_{-\infty}^L(M_a-1)d\lambda
\]
for each of the three rays, since the extended profile is zero above
its cutoff. The negative-tail integrals are finite by finite energy;
one never subtracts divergent integrals. Subtracting gives
\[
 \int_{\mathbf R}(M_{s_j}-M_{r_j})=E_\alpha(\lambda_j)-E_\alpha(\eta_j),
 \quad
 \int_{\mathbf R}(M_{s_j}-M_r)=E_\alpha(\lambda_j)-E_\alpha(\phi).
\]
The pointwise inequality
$|M_{r_j}-M_r|\leq(M_{s_j}-M_{r_j})+(M_{s_j}-M_r)$ proves
\eqref{eq:mass-bound}.

For nonincreasing functions $f_j,f$ with $\|f_j-f\|_{L^1}\to0$,
at a continuity point $x$ the values converge: the average over
$(x,x+h)$ is at most $f_j(x)$, which is at most the average over
$(x-h,x)$. Pass to the limit in the averages, then let $h\downarrow0$.
L gives continuity and positive mass at all interior levels. Uniform
convergence on a compact interval follows by a finite grid and
monotonicity, using uniform continuity of $M_r$. Also
$\sup\eta_j\to\sup\phi$, so every such compact interval is below
the approximants' cutoffs for large $j$.
\end{proof}

\begin{corollary}\label{cor:curves}
In every dimension the same sequence converges at every interior
level in $d_{\mathcal S}$, and all its mixed level masses converge:
\[
 d_{\mathcal S}([\chi_{r_j}(\lambda)],[\chi_r(\lambda)])\to0
 \quad(\lambda<\sup\phi).
\]
The convergence is uniform on compact interior intervals. This is
convergence in a topology detecting $P$-equivalence at positive mass,
not exact $P$-equivalence at each index or a zero-mass endpoint claim.
\end{corollary}
\begin{proof}
First consider any ordered pair $u\leq v$ of $\omega$-psh functions
with $M_u>0$, and put $d=M_v-M_u\geq0$ and
$M_k(u)=V^{-1}\int\omega^{n-k}\wedge
\langle\omega_u^k\rangle$. If $b>1$ and $b^n d<M_v$,
the just-stated Lemma 4.3 applies. With
$h=P_\omega(bu+(1-b)v)$ one has
$b^{-1}h+(1-b^{-1})v\leq u$. Mixed-mass monotonicity,
multilinearity and positivity give, for $0\leq k\leq n$,
\[
 M_k(v)\geq M_k(u)\geq(1-b^{-1})^kM_k(v),\qquad
 0\leq M_k(v)-M_k(u)\leq k/b.
\]
The last bound uses $M_k(v)\leq1$. For $0<d<M_v/2$ choose
$b=(M_v/(2d))^{1/n}>1$. Equation \eqref{eq:ordered-dS} then gives
\[
 d_{\mathcal S}([u],[v])\leq\frac{n}{2}
                   \left(\frac{2d}{M_v}\right)^{1/n}.
\]
If $d=0$, every $b>1$ is allowed, and letting $b\to\infty$
proves that all mixed-mass differences and $d_{\mathcal S}$ vanish.

Now apply this estimate to each of $\chi_{r_j}(\lambda)$ and
$\chi_r(\lambda)$ below the common upper level
$\sigma_j(\lambda)$, and use triangle inequality.
The proof of Theorem \ref{thm:profiles} also gives
$\|M_{s_j}-M_r\|_{L^1}\to0$. The same monotone-function argument
therefore makes all three masses converge at each interior level,
uniformly on compact intervals. On every such interval their
limiting mass has a positive minimum. Thus the two ordered top-mass
defects tend uniformly to zero with a uniform positive denominator
in the bound above. This proves the $d_{\mathcal S}$ claim.
Applying the mixed-mass bound to both ordered pairs and subtracting
through their common upper level proves uniform convergence of
every $M_k$ as well. In dimension one the sharper bound is simply
$d_{\mathcal S}([u],[v])=(M_v-M_u)/2$ for the ordered pairs.
\end{proof}

\begin{proposition}\label{prop:initial}
For arbitrary genuine $w_j\to w$ in analytic $d_1$ and
$\eta_j\to\phi$ strongly non-Archimedean, the actual rays satisfy
\[
 d_1(r(w_j,\eta_j)_t,r(w,\phi)_t)\leq A_j+tB_j,
\]
where $B_j$ is as above and
$A_j=2E(\max(w_j,w))-E(w_j)-E(w)\to0$.
\end{proposition}
\begin{proof}
The common upper maximal ray starts at $a_j=\max(w_j,w)$ and has
datum $\lambda_j$. Ordered energies give the displayed inequality.
The analytic supporting inequality and $I_1\leq Cd_1$ give
\[
 0\leq E(a_j)-E(w)\leq\int|w_j-w|\MA(w)\leq Cd_1(w_j,w).
\]
Energy continuity then gives $A_j\to0$.
\end{proof}
This proposition proves uniform $d_1$ convergence on bounded time
intervals also for nonsmooth initial points. It proves no relative
level-envelope or entropy theorem for such initial points by itself.

\section{Discrepancy integrability retained by entropy recovery}
\begin{lemma}\label{lem:UI}
If probability measures $\nu_j\rightharpoonup\nu$ on a compact
space and a nonnegative lower semicontinuous function $A$ satisfy
$\int A\,d\nu_j\to\int A\,d\nu<\infty$, then
\[
 \lim_{R\to\infty}\sup_j\int_{\{A>R\}}A\,d\nu_j=0.
\]
For $\nu_j=\MA_\alpha(\eta_j)$ in R this applies to $A=A_X$.
\end{lemma}
\begin{proof}
For $M>0$, $A\wedge M$ is bounded lower semicontinuous. Weak
convergence and its integral lower semicontinuity give
\[
 \limsup_j\int(A-M)_+\,d\nu_j
 \leq\int A\,d\nu-\int(A\wedge M)\,d\nu\to0
 \quad(M\to\infty).
\]
On $\{A>2M\}$, $A\leq2(A-M)_+$. This controls every sufficiently
late index; each of the finitely many earlier measures has a finite
$A$ integral and its own tail tends to zero. First choose $M$ for
the limiting integral and late indices, then enlarge the cutoff to
handle the finite early set. This proves the uniform assertion.
On $X^\beth$ the measures are regular probability measures and
$A_X$ is the supremum of continuous model tests, so this lower
semicontinuity/Portmanteau application is valid. Infinite values are
allowed, but have measure zero for these finite integrals.
\end{proof}
This is uniform integrability of discrepancy against changing
non-Archimedean measures. It is not a statement that analytic
densities converge in measure or that their entropy densities are
uniformly integrable at all finite times.

\section{Uniform mixed control and a joint entropy limit}\label{sec:mixed}
For this section all rays start at zero in a fixed K\"ahler background.
For a smooth initial $w$ use the background change specified above.
\begin{lemma}\label{lem:uniform-mixed}
With $r_j,r,B_j$ as in Theorem \ref{thm:profiles}, $t\geq1$, and
fixed smooth contracted expression $G$ in \eqref{eq:M},
\begin{equation}\label{eq:G-uniform}
 \frac{|G(r_{j,t})-G(r_t)|}{t}
 \leq C\left(B_j^q+|\sup\eta_j-\sup\phi|
                    +|E_\alpha(\eta_j)-E_\alpha(\phi)|+t^{-1}\right).
\end{equation}
For every fixed smooth-model divisor test $C$ and compatible $\Psi$,
\begin{equation}\label{eq:Psi-uniform}
 \frac1t\left|\int\Psi_t\MA(r_{j,t})-\int\Psi_t\MA(r_t)\right|
          \leq C_C B_j^q.
\end{equation}
The constants are independent of the changing models of $\eta_j$.
\end{lemma}
\begin{proof}
By supremum affinity, finite-energy affinity and the compact psh
supremum-minus-mean bound $T_\omega$,
\[
 J(r_{j,t}),J(r_t),J(s_{j,t})\leq C t,\qquad
 |\bar r_{j,t}-\bar r_t|
 \leq t|\sup\eta_j-\sup\phi|+2T_\omega,
\]
where bars denote averages against $\mu_0$. Constants are uniform
because the non-Archimedean energies and suprema are bounded.
The mixed-pairing estimate Q applies after subtracting these means.
The slot energy distance is at most a fixed multiple of
$d_1(r_{j,t},r_t)\leq tB_j$ (the $J$ distance is constant invariant).
Thus its normalized mixed error is at most
$C(tB_j)^q(t+T_\omega)^{1-q}\leq CtB_j^q$ for $t\geq1$.
Replace the $n$ varying slots one at a time.

For $T_\xi$ the restoration of means changes the pairing by fixed
cohomological constants times the mean difference; splitting a
signed $\xi$ into two positive smooth forms is legitimate.
$E(r_{j,t})-E(r_t)=t(E_\alpha(\eta_j)-E_\alpha(\phi))$ exactly.
These statements prove \eqref{eq:G-uniform}.

For the second assertion the expression is the Deligne pairing
$(0,\Psi_t)(\omega,r_t)^n/V$. A constant change in any of the
$r_t$ slots contributes an integral of $\ddc\Psi_t$ against a closed
product, which is zero on the compact noncentral fiber. A constant
change in the test slot contributes the same number to both
probability measures and cancels in the difference. Consequently no
mean-restoration error remains in \eqref{eq:Psi-uniform}.
Represent the fixed signed test as a difference of smoothly
compatible positive model weights on one common model, as in the
proof of \cite{MP}, Theorem 5.3.4. These fixed slots have $J=O(t+1)$.
Their forms and models are held fixed before $j$ or $t$ varies, so Q
has the asserted uniform constant.
\end{proof}

Here is the normalization of the slopes in Q, needed also for the
entropy argument. For smooth compatible model weights the mixed
Deligne pairing has base $\ddc$ equal to the pushforward of the product
of the total smooth class representatives. Stokes' theorem on the
base minus a disk, with metrics constant near infinity, gives slope
equal to the intersection. Indeed $\ddc F=(4\pi)^{-1}\Delta F\,dA$
and the inner boundary flux is $-rF_r/2=F_t$ for
$t=-\log r^2$. There is no extra factor two. For singular maximal
slots use decreasing model data and decreasing initial points;
ordered $d_1$ is an initial energy error plus $t$ times the datum
energy error. Q bounds the mixed error by that distance to power
$q$ times an $O(t+1)$ quantity to power $1-q$. Divide by $t$,
take the time limit first and then the decreasing approximation
limit. This reconstructs the singular slope passage in \cite{MP}.
In MP's $i\partial\bar\partial/\pi$, logarithm and
$s=-\log|\tau|$ convention use background $2\omega$ and potential
$r_{2s}$. Its datum is $2\phi$ in class $2\alpha$. Currents double,
probability measures agree, and division by $t=2s$ cancels the
doubled mixed slope. This transport concerns mixed pairings only.

\begin{lemma}\label{lem:log-volume}
On a fixed smooth dominating SNC model $\mathcal Y$, put
$B=\mathcal Y_0$ and
$C_{\mathcal Y}=K_{\mathcal Y/(X\times\mathbf P^1)}+B_{\rm red}-B$.
There is a smooth positive volume $\mu_{\mathcal Y,t}$ on every
noncentral fiber with
\[
 \Psi_t=\log(\mu_{\mathcal Y,t}/\mu_0),\quad
 \Psi+\log|s_{C_{\mathcal Y}}|^2\text{ smooth},\quad
 \int\mu_{\mathcal Y,t}\leq C(1+t)^n.
\]
If $\Omega$ is a fixed smooth positive form on the model, then
$\Omega_t^n\leq C'\mu_{\mathcal Y,t}$ for sufficiently large $t$.
The associated PL test is $f_{C_{\mathcal Y}}=A_X\circ r_{\mathcal Y}$.
\end{lemma}
\begin{proof}
Equip $K_{\mathcal Y/\mathbf P^1}+B_{\rm red}-B$ with a smooth
Hermitian metric, defining its fiber volume. Near an SNC stratum
write $\tau=z_0^{b_0}\cdots z_p^{b_p}$, absorbing a unit in one
coordinate. On the fiber its log-canonical frame is, up to a nonzero
constant,
\[
 \frac{dz_1}{z_1}\wedge\cdots\wedge\frac{dz_p}{z_p}
           \wedge dz_{p+1}\wedge\cdots\wedge dz_n.
\]
The base fiber volume has coefficient $s_{C_{\mathcal Y}}$ against
this frame, of order $a_i+1-b_i$ on a component of canonical order
$a_i$ and multiplicity $b_i$. This proves compatibility and the
normalized Gauss discrepancy $(a_i+1-b_i)/b_i$ over $X$.

In bounded polydisks the independent radial variables
$y_i=-\log|z_i|^2$ occupy, after fixed shifts, the weighted simplex
$\sum_{i=1}^p b_i y_i\leq t+O(1)$, $y_i\geq0$. Logarithmic
radial volume is a constant times $dy_1\cdots dy_p$. Angular and
free-coordinate factors are bounded and there are only finitely many
branches for the dependent $z_0$. The integral is $O((1+t)^p)$,
$p\leq n$. A finite cover proves the asserted volume bound,
retaining all central multiplicities.
On each fiber
$dz_0=-z_0\sum_{i=1}^p(b_i/b_0)dz_i/z_i$.
Every ambient coordinate $n$-wedge has a bounded coefficient against
the logarithmic frame (the other coefficients contain products of
the bounded $z_i$). A smooth ambient positive form has bounded
coefficients in these charts, proving its fiber-volume domination.
\end{proof}

\begin{proposition}\label{prop:model-slope}
For an actual smooth relatively K\"ahler model datum $\eta$ and
its ray $r(0,\eta)$,
$\lim_{t\to\infty}H(r_t)/t=H_\alpha(\eta)$.
\end{proposition}
\begin{proof}
Resolve the central support and pull back the ray. The regularity
input SD implies a positive residual current bounded by $C\Omega$
near that fiber, and hence on noncentral fibers
$\MA(r_t)\leq C_1\Omega_t^n\leq C_2\mu_{\mathcal Y,t}$.
For continuous bounded-Hessian potentials the determinant comparison
follows by local convolution, its almost-everywhere determinant and
Bedford--Taylor continuity; the continuous current inequalities
restrict to each noncentral fiber.
Put $g_t=\MA(r_t)/\mu_{\mathcal Y,t}\leq C_2$. The entropy
decomposition gives
\[
 H(r_t)=\int\log g_t\MA(r_t)+\int\Psi_t\MA(r_t)
       \leq\log C_2+\int\Psi_t\MA(r_t).
\]
For each fixed $t$ the volumes are smooth and positive and
$g\log g$ has integrable negative part; this is a valid finite
entropy decomposition. The calibrated slope gives the upper bound
by $\int f_{C_{\mathcal Y}}\MA_\alpha(\eta)=H_\alpha(\eta)$,
since the model measure is supported on these components.
For the lower bound Gibbs' entropy inequality gives
\[
 H(r_t)\geq\int\Psi_t\MA(r_t)
              -\log\int\mu_{\mathcal Y,t}.
\]
The last term is $O(\log(1+t))$. This proves the same lower slope.
\end{proof}

\begin{theorem}\label{thm:joint-entropy}
Let $\phi\in\cE^{1,\NA}(\alpha)$ have finite discrepancy entropy,
and let the actual model data $\eta_j$ satisfy R. Their actual rays
from a fixed smooth K\"ahler initial potential $w$ obey, for every
choice $j_\ell\to\infty$, $t_\ell\to\infty$,
\[
 \frac{H(r(w,\eta_{j_\ell})_{t_\ell})}{t_\ell}
          \longrightarrow H_\alpha(\phi).
\]
The entropy slope of $r(w,\phi)$ has this value as well.
This statement concerns rescaled entropies on one fixed manifold,
not continuity of entropy at each fixed time or a hybrid measure
theorem for degenerating families.
\end{theorem}
\begin{proof}
First $w=0$. Let $h_j=H_\alpha(\eta_j)$ and
$g_j=G_\alpha(\eta_j)$. Proposition \ref{prop:model-slope}, the
calibrated contracted slopes, and convexity of $M$ imply
\[
 M(r_{j,t})\leq M(0)+t(h_j+g_j).
\]
This uses only finite-index model regularity; the regularity constants
may depend on $j$. Strong continuity of the fixed contracted pairings
gives $g_j\to G_\alpha(\phi)$ and R gives $h_j\to H_\alpha(\phi)$.
Lemma \ref{lem:uniform-mixed} and the target ray's contracted slope
give, along every such joint sequence,
$G(r_{j_\ell,t_\ell})/t_\ell\to G_\alpha(\phi)$. Thus the joint
limsup of $H/t$ is at most $H_\alpha(\phi)$.

For the liminf fix one SNC model $\mathcal Y$, before either limit.
Gibbs' inequality and Lemma \ref{lem:log-volume} give
\[
 \frac{H(r_{j,t})}{t}\geq\frac1t\int\Psi_t\MA(r_{j,t})
     -\frac{\log C+n\log(1+t)}t.
\]
Lemma \ref{lem:uniform-mixed} replaces its integral by the target
ray's with error at most $C_{\mathcal Y}B_j^q$. The calibrated target
slope therefore proves a joint lower bound by
$\int(A_X\circ r_{\mathcal Y})\MA_\alpha(\phi)$.
Take the supremum over fixed dominating SNC models only after this
joint liminf. The directed discrepancy-integral identity in F gives
$H_\alpha(\phi)$. This proves the joint limit without exchanging
varying models and time in the log-volume constants.

For the target ray itself, the same Gibbs argument gives its lower
slope. At each fixed $t$ pass to the model approximation in the
displayed Mabuchi inequality using \eqref{eq:ray-bound} and lower
semicontinuity, obtaining
$M(r_t)\leq M(0)+t(H_\alpha(\phi)+G_\alpha(\phi))$.
Divide by $t$ and subtract the contracted slope for the upper bound.
This also shows each of its finite-time entropies is finite.

For fixed smooth $w$ pass to $\omega_w$ and potentials $r_t-w$.
The same actual configurations, usual data, probability energy slopes
and non-Archimedean entropy apply, since the class is unchanged and a
smooth $w$ has no divisorial coefficient. The two smooth reference
probabilities have a bounded logarithmic ratio, so their entropies
differ by a uniformly bounded term, which vanishes after division
by $t_\ell$. This completes the stated initial-point transport.
\end{proof}
No analytic entropy uniform-integrability or density-in-measure
assumption is hidden here. The upper bound uses convexity and the
recovered slopes; the lower bound uses fixed model tests and uniform
mixed control. Lemma \ref{lem:UI} supplies the separate
non-Archimedean discrepancy integrability. Infinite $H_\alpha$ and
arbitrary singular initial potentials are outside this joint theorem.

\begin{corollary}[Simultaneous recovery at interior levels]\label{cor:combined}
On a fixed smooth compact K\"ahler $(X,\omega)$, fix a smooth
K\"ahler initial potential $w$ and
$\phi\in\cE^{1,\NA}([\omega])$ with finite discrepancy entropy.
There are actual smooth relatively K\"ahler model configurations
$\eta_j$ whose actual rays all start at this same $w$, with strong
energy and entropy recovery. For their inverse Legendre levels
relative to $w$, every level strictly below $\sup\phi$ converges
in $d_{\mathcal S}$ and all mixed raw masses converge, uniformly
on compact interior level intervals. Their full top-mass profiles
converge in $L^1$ of the level parameter. Their discrepancy tails
are uniformly integrable. Every joint limit of their rescaled
analytic entropy as model index and time tend to infinity equals
$H_\alpha(\phi)$.

This is a simultaneous theorem in the stated interior-level topology,
not a complete solution of the original P19.2 endpoint and prescribed
singularity programme.
\end{corollary}
\begin{proof}
Choose the actual configurations of R. Apply Theorem
\ref{thm:profiles}, Corollary \ref{cor:curves} and Lemma
\ref{lem:UI}. For the level assertions at the genuine smooth initial
point use $\omega_w$ and $r_t-w$, with reference volume
$\omega_w^n/V$; this changes neither the class nor the actual models
or datum and gives exactly the relative levels after adding $w$.
The joint entropy conclusion is Theorem \ref{thm:joint-entropy},
with the bounded reference-volume change proved there.
No approximation of the initial potential is used in this signature.
\end{proof}

\section{An actual positive-discrepancy model and its endpoint}\label{sec:explicit}
\begin{theorem}\label{thm:explicit}
On $\mathbf P^1$, let $\omega=\ddc\log(1+|z|^2)$, of mass one,
and let $c\in\mathbf Q\cap(0,1)$. There is an actual relatively
ample model ray from zero on the smooth blowup of $(z,\tau)=(0,0)$,
with vertical divisor $D=-cF$. It satisfies
\[
 E(u_t)=-\tfrac12c^2t,\qquad H^\infty(u)=c,
 \quad E_\alpha(f_D)=-\tfrac12c^2,\quad H_\alpha(f_D)=c.
\]
Its Legendre curve is $0$ for $\lambda\leq-c$, $w_{c+\lambda}$
for $-c<\lambda\leq0$, and $-\infty$ for $\lambda>0$, where
\[
 w_b(z)=\begin{cases}
 b\log|z|^2-F(b)-\log(1+|z|^2),&|z|^2\leq b/(1-b),\\
 0,&|z|^2\geq b/(1-b),
 \end{cases}
\]
\[
 F(b)=b\log b+(1-b)\log(1-b).
\]
In particular the endpoint at $\lambda=0$ has positive raw mass
$1-c$ and a logarithm-square pole of coefficient $c$ at zero.
\end{theorem}
\begin{proof}
Put $\rho=\log|z|^2$ and
\[
 q_t(\rho)=\sup_{0\leq x\leq1}
       \{x\rho-F(x)-t(c-x)_+\},\qquad
 u_t=q_t(\rho)-\log(1+e^\rho).
\]
At the endpoints of $[0,1]$ use $0\log0=0$. Since
$F'(x)=\log(x/(1-x))$, with $A=c/(1-c)$ and $y=|z|^2$,
\[
 q_t(\log y)=\begin{cases}
 \log(1+ye^t)-ct,&y<Ae^{-t},\\
 c\log y-F(c),&Ae^{-t}\leq y\leq A,\\
 \log(1+y),&y>A.
 \end{cases}
\]
These pieces match with their first derivatives. $u_0=0$ and
$-ct\leq u_t\leq0$. The fiber curvature is positive, with radial
probability measure $d x$, where $x=\partial_\rho q_t$.
Its mass is $c$ on the inner disk, zero on the middle annulus and
$1-c$ outside. No circle mass occurs because the first derivatives
match. The functions extend across zero and infinity at each finite
time and are bounded there.

The total local weight is $q_{-\log|\tau|^2}(\log|z|^2)$.
On its inner region it is
$\log(|\tau|^2+|z|^2)-(1-c)\log|\tau|^2$;
on the middle it is $c\log|z|^2-F(c)$ and on the exterior it is
$\log(1+|z|^2)$. Each has semipositive complex Hessian and rank at
most one on the punctured total space. The matching first derivatives
ensure no interface current. The potentials are locally
$C^{1,1}$ on any compact noncentral annulus, so their
Bedford--Taylor square is the almost-everywhere determinant and is
zero. The actual weak-ray conclusion also follows directly from the
$P$-model level calculation below and the analytic Legendre
correspondence L; it does not depend on treating an arbitrary
subgeodesic as a weak ray.

In the blowup chart $z=s\tau$, subtracting $c\log|\tau|^2$
from $q$ gives $\log(1+|s|^2)$ for $|s|^2<A$ and
$c\log|s|^2-F(c)$ on the adjacent region. These match smoothly to
first order and are bounded. In the other blowup chart, subtraction
of $c\log|z|^2$ gives the same bounded compatibility; away from the
exceptional divisor there is no divisorial coefficient. Thus the
actual datum is $f_{-cF}$, with values $-c$ on the exceptional Gauss
valuation $\ord_0$ and $0$ on the strict central component.

For rational $c$, a common denominator makes
$p_1^*\mathcal O(1)-cF$ a genuine relative $\mathbf Q$-line bundle.
It has degrees $c$ and $1-c$ on the two central components and
degree one on every other fiber. Both central multiplicities are one.
Each noncentral fiber is a projective line. The central fiber is two
projective lines meeting in one node. Positive integral component
degrees make a sufficiently high power very ample on this nodal
curve: sections on either component vanishing at the node extend by
zero to the other, separate points and tangent directions away from
the node, and separate the two tangent directions at the node;
sections with a shared nonzero node value generate there.
Thus the line bundle is ample on every fiber. The proper-morphism
fiber criterion \cite{Stacks} gives relative ampleness locally on the
base, and hence relative ampleness. Adding a sufficiently positive
base class yields a K\"ahler
class on the compact blowup. This proves actual admissibility and
positivity; $c=1$ would only be the boundary nef case and is not used.
The defining model measure is
$(1-c)\delta_{v_{\rm triv}}+c\delta_{\ord_0}$, so its entropy is
$c$ because $A_X(\ord_0)=1$ and $A_X(v_{\rm triv})=0$.

Energy differentiation is legitimate for these locally bounded
piecewise smooth weights: $\partial_tu_t=-(c-x)_+$ at the maximizing
$x$. Hence $E'(u_t)=-\int_0^c(c-x)dx=-c^2/2$ and $E(u_0)=0$.
It also follows from the Legendre mass formula below. The model
intersection energy is $D^2/2=-c^2/2$, since $F^2=-1$,
$p_1^*\alpha\cdot F=0$ and $(p_1^*\alpha)^2=0$ on the compact
model. These give the same normalization independently.

The density relative to $\omega$ is one on the exterior, zero on
the middle annulus and
$e^t(1+y)^2/(1+e^t y)^2$ on $0\leq y<Ae^{-t}$.
Writing $k=e^ty$ in the inner region gives
\[
 H(u_t)=ct+2\int_0^A
       \frac{\log(1+e^{-t}k)-\log(1+k)}{(1+k)^2}dk.
\]
The integral is bounded uniformly in $t\geq0$ and every finite-time
entropy is finite. Thus $H(u_t)/t\to c$, an actual positive
discrepancy calibration.

To compute the Legendre transform, for $-c<\lambda<0$ put
$b=c+\lambda$. On the inner region the critical time is
$t=\log(b/(1-b)) -\log y$, if this is nonnegative.
Substitution gives $w_b$; if that time is negative, the infimum is at
$t=0$ and is zero. The endpoint $\lambda=0$ is the decreasing
$t\to\infty$ limit $w_c$. For $\lambda\leq-c$ the infimum is zero
at $t=0$, and for $\lambda>0$ it is $-\infty$ since $u_t\leq0$.
The curvature of $w_b$ is the atom $b\delta_0$ plus $\omega$ on
$y>b/(1-b)$, with no interface mass. Removing the pole removes
the atom, leaving raw non-pluripolar mass $1-b$.
Here is a direct check of their $P$-model property, without assuming
the ray conclusion. For $0<b<1$ let $B=b/(1-b)$. A rooftop
competitor $h\leq0$, $h\leq w_b+C$ has local subharmonic weight
$Q=h+\log(1+|z|^2)$. On $|z|^2<B$, $Q-b\log|z|^2$ is
subharmonic: subtract the atom $b\delta_0$, which is present with
at least this coefficient by the singularity bound. It is bounded
above near zero, and thus extends subharmonically across zero.
On the boundary $Q\leq\log(1+B)$. The maximum principle gives
$Q\leq b\log|z|^2+\log(1+B)-b\log B
=b\log|z|^2-F(b)$ throughout the disk. Outside it $h\leq0=w_b$.
Every competitor therefore lies below $w_b$, and $w_b$ itself is a
competitor. Hence $P_\omega[w_b]=w_b$, including $b=c$.

The calculated Legendre curve is decreasing, parameter-concave and
upper semicontinuous, as a Legendre infimum of the displayed convex
functions of time. Its nontrivial levels are the psh $w_b$; it is
zero for $\lambda\leq-c$, has endpoint $w_c$ and is $-\infty$
above zero. It is consequently a bounded $P$-maximal test curve.
Pointwise convex duality, or substitution into the formulas, gives
its inverse Legendre transform equal to $u_t$. Correspondence L
proves that this is the actual bounded weak ray from zero.
This calculation retains the endpoint and poles as well as masses.
\end{proof}

\begin{corollary}\label{cor:endpoint}
There is a sequence of actual smooth relatively ample models with
strong convergence and exact entropy recovery whose fixed target
endpoint $P$-envelope and raw mass do not converge, even when that
target endpoint has positive mass and positive discrepancy data.
\end{corollary}
\begin{proof}
On the model of Theorem \ref{thm:explicit}, let $B=\mathcal Y_0$ and
$D_j=D-B/j$. A full central-fiber translation leaves relative
ampleness unchanged. Its datum is $\eta_j=\eta-1/j$, its actual
initial-zero ray is $u_{j,t}=u_t-t/j$, and
\[
 E_\alpha(\eta_j)=E_\alpha(\eta)-1/j,\quad
 \MA_\alpha(\eta_j)=\MA_\alpha(\eta),\quad
 H_\alpha(\eta_j)=c.
\]
These are actual rational algebraic configurations on the same smooth
model; uniform constant convergence implies strong convergence.
The analytic measures and entropies are exactly those of $u_t$ at
every time. But $\chi_j(\lambda)=\chi(\lambda+1/j)$, so
$\chi_j(0)\equiv-\infty$, $M_j(0)=0$, while
$\chi(0)=w_c$ and $M(0)=1-c>0$. At the moving cutoff
$\lambda_j=-1/j$, the endpoint is exactly $w_c$ and its mass is
exactly $1-c$. Fixed-level and moving-level signatures differ.
\end{proof}
This corollary refutes an implication from strong energy--entropy
recovery to fixed endpoint preservation. It is not a nonexistence
theorem for a carefully chosen approximation: the constant sequence
$\eta_j=\eta$ preserves all these data exactly. It therefore does
not resolve the original P19.2 by counterexample.

\section{What the polar Cantor energy defect obstructs}
\begin{proposition}\label{prop:defect}
Let $u$ be an actual finite-energy weak ray from $w$, with usual
datum $\phi\in\cE^{1,\NA}(\alpha)$, and let
$\widehat u=r(w,\phi)\geq u$. Write
$c_u=E(u_1)-E(u_0)$ and
$\delta=E_\alpha(\phi)-c_u\geq0$.
If maximal model rays $r_j=r(w_j,\eta_j)$ have
$w_j\to w$ in $d_1$ and $\eta_j\to\phi$ strongly, then
\[
 \lim_{t\to\infty}\frac{d_1(\widehat u_t,u_t)}t=\delta,
 \quad
 \left|\liminf_{t\to\infty}\frac{d_1(r_{j,t},u_t)}t-\delta\right|\leq B_j,
\]
and the same bound holds for the limsup in $t$.
In particular both quantities tend to $\delta$ as $j\to\infty$.
If $\delta>0$, such rays cannot preserve the true ray energy slope
while recovering that usual datum strongly, and cannot approximate
$u$ in chordal distance.
\end{proposition}
\begin{proof}
Ordered energy and affinity give
$d_1(\widehat u_t,u_t)=t\delta$. Proposition \ref{prop:initial}
gives $d_1(r_{j,t},\widehat u_t)\leq A_j+tB_j$.
The reverse and direct triangle inequalities sandwich
$d_1(r_{j,t},u_t)/t$ between
$\delta-B_j-A_j/t$ and $\delta+B_j+A_j/t$.
Take liminf and limsup in $t$, then $j\to\infty$.
Also every maximal $r_j$ has energy slope $E_\alpha(\eta_j)$,
which converges to $E_\alpha(\phi)$ rather than $c_u$ if
$\delta>0$. This proves nonexistence for this precisely specified
mechanism, not for all possible admissible approximation signatures.
\end{proof}

The frozen input supplies a complete polar Cantor construction on
$\mathbf P^1\times Y$, with $Y$ any compact K\"ahler manifold, and
$a\in(0,1)$. Its old convention is
$\ddc_L=i\partial\bar\partial/\pi$, $s=-\log|\tau|$,
mass-one curve background $\omega$ and product background $\Theta$.
Let its old potentials be $v_s$. In our convention keep the same
forms and define
\[
 u_t=2v_{t/2},\qquad \psi_{\rm new}=2\psi_{\rm old}.
\]
Then $\Theta+\ddc u_t=\Theta+\ddc_Lv_{t/2}$; total currents
and the homogeneous equation agree. Probability measures agree,
$E_\Theta(2v)=2E_{\Theta,L}(v)$, and consequently
\[
 -at\leq u_t\leq0,\quad E_\Theta(u_t)=-a^2t/2,\quad
 u^{\NA}=0,\quad H(u_t)/t\to+\infty.
\]
Generic coefficients of $2U_L$ measured against $\log|f|^2$ equal
those of $U_L$ against $\log|f|$, so its actual usual datum remains
zero. Its Legendre levels become $2$ times the old levels with the
same level parameter; the prescribed positive-mass polar envelopes
are not reassigned or discarded. Its maximalization from zero is
the zero ray and $\delta=a^2/2$.

Proposition \ref{prop:defect} therefore gives the sharp limiting
distance $a^2/2$ for model maximal rays recovering this usual datum.
This use of the Cantor example depends on its supplied analytic
construction, not on a historical PASS. Its pole measure is atom
free, all exponential powers of the polar potential are integrable,
and its entropy slope is infinite, as proved in that supplied text.
In particular it is not a counterexample in a regime requiring finite
true ray entropy, nor is it a ray maximal for a log-discrepancy datum.
The original residual target concerns prescribed model and singularity
data; the excerpt does not assert that all nonmaximal weak rays with
infinite entropy must be approximated by maximal models having only
their ordinary datum. Thus the proposition rules out an overly broad
mechanism and does not close P19.2. Assigning another datum to this
ray would change, rather than preserve, its original signature.

\section{Prescribed singularities, capacity, and the remaining gaps}
For log-discrepancy-model inputs R retains the actual model
configurations, their usual datum, energy, probability entropy, and
discrepancy tails. Theorem \ref{thm:profiles} additionally retains
the limiting entire raw mass profile, including every interior level
mass, and Corollary \ref{cor:curves} the interior $d_{\mathcal S}$
singularity topology and all mixed level masses in every dimension.
This uses a quantitative ordered-envelope argument, rather than
equating top mass with the full metric formula.
Our proof does not assert exact envelopes or exact poles at each
index. No replacement of $P$-equivalence by
$\mathcal I$-equivalence is made. Equal multiplier ideals alone
cannot identify general positive-mass $P$-envelopes.

The relevant fresh moving-singularity result is Pang--Sun--Wang--Zhou
\cite{PSWZ}, Theorem 1.1: on a fixed compact K\"ahler manifold,
smooth big forms $\theta_j\to\theta$ with
$-\varepsilon_j\omega\leq\theta_j-\theta\leq\varepsilon_j\omega$,
normalized positive-mass model potentials $\varphi_j,\varphi$, and
non-pluripolar positive measures $\mu_j,\mu$ converging in total
variation with masses respectively the prescribed raw masses,
the normalized relative full-mass solutions converge in capacity
if and only if the model potentials do. Its applicability requires
the total variation hypothesis and the prescribed-capacity
signature. Theorem \ref{thm:profiles} provides neither general
level-capacity convergence nor total variation convergence of the
corresponding right-hand sides. Accordingly this theorem is not
applied to infer either missing conclusion. The same preprint's
arbitrary-mass envelope idempotence is not an entropy recovery or
endpoint-continuity theorem. We rely only on established
positive-mass idempotence where it enters our proofs.

For clarity the unsolved requirements are concrete. They include a
choice and justification of the original residual relative
$P$-singularity topology for the entire original signature;
preservation of any further prescribed envelopes and exact poles;
the endpoints
and zero-mass levels; and genuine singular initial points in these
level and entropy assertions. Exact indexwise preservation has not
been proved even where our convergence holds. The joint entropy
theorem divides entropy by time and assumes a fixed smooth start;
it does not prove a fixed-time hybrid entropy limit or prescribed
relative-level entropy continuity. The original target's full
simultaneous construction or a counterexample to its existence
is still missing. No genuinely degenerating-family conclusion
(P19.3), automorphism-relative or singular stability extension
(P19.4), or new completion of P19.1 is claimed.

\section{Verification and reproducibility}
This is an unverified research draft. Analytic proofs above, exact
external premises and the supplied recovery theorem establish the
claimed partial signature; execution success is not a proof
certificate. The explicit calibration and endpoint formulas are
checked at moderate finite parameters by a small Python program.
These checks are numerical observations, not exhaustive finite
results or replacements for the analytic proofs.

The working directory contains byte-identical \texttt{input.json}
and \texttt{input-original.json}, decoded
\texttt{assignment-original.txt}, and \texttt{task-original.json}.
The supplied mathematical units are extracted locally without
access to another campaign. \texttt{source-application-ledger.json}
states which primary proof units were actually read and which
foundations are imported; no recursive audit of every cited
foundation is claimed. Shell PDF downloads were blocked; saved
web-tool text excerpts, not inaccessible PDF screenshots, are the
fresh public evidence. The source tool requests can be repeated
from the ledger with the existing web tool. No resource-heavy
compute facility was exposed or needed, and none was launched.

Run \texttt{python check\_models.py} and
\texttt{python validate\_deliverables.py}. Compile with
\texttt{pdflatex -interaction=nonstopmode -halt-on-error P19.2-paper.tex}
twice. The build script retains both runs' output and logs.
\texttt{P19.2-argument.md} is the readable companion;
\texttt{status.json} records PARTIAL with the exact gaps.

\begin{thebibliography}{99}
\bibitem{MP} P. Mesquita--Piccione, \emph{A non-Archimedean theory
of complex spaces and the cscK problem},
\href{https://arxiv.org/pdf/2409.06221v2}{arXiv:2409.06221v2}.
\bibitem{MPWN} P. Mesquita--Piccione and D. Witt Nystr\"om,
\emph{A transcendental non-Archimedean Calabi--Yau theorem with
applications to the cscK problem},
\href{https://arxiv.org/pdf/2509.09442v1}{arXiv:2509.09442v1}.
\bibitem{T} A. Trusiani, \emph{On the transcendental
Yau--Tian--Donaldson conjecture},
\href{https://arxiv.org/pdf/2609.31089v1}{arXiv:2609.31089v1},
25 September 2026. Mathematical construction used through the
precise supplied theorem R, including its local discrepancy repair.
\bibitem{SD} Z. Sj\"ostr\"om Dyrefelt,
\emph{K-semistability of cscK manifolds with transcendental
cohomology class},
\href{https://arxiv.org/pdf/1601.07659v2}{arXiv:1601.07659v2},
Lemma 4.6 and Remark 4.7.
\bibitem{BDL} R. Berman, T. Darvas and C. H. Lu,
\emph{Convexity of the extended K-energy and the large time behavior
of the weak Calabi flow}, Geom. Topol. 21 (2017), 2945--2988,
\href{https://msp.org/gt/2017/21-5/gt-v21-n5-p08-s.pdf}{primary text}.
\bibitem{D} T. Darvas, \emph{The Mabuchi geometry of finite energy
classes}, \href{https://arxiv.org/pdf/1409.2072v2}{arXiv:1409.2072v2},
Theorem 5.5 and energy/ordered metric identities.
\bibitem{DX} T. Darvas and M. Xia, \emph{The closures of test
configurations and algebraic singularity types}, Adv. Math. 397
(2022), 108198,
\href{https://research.chalmers.se/publication/528300/file/528300_Fulltext.pdf}
{published primary text}, analytic proof units in Section 3.1.
\bibitem{DXZ} T. Darvas, M. Xia and K. Zhang,
\emph{A transcendental approach to non-Archimedean metrics of
pseudoeffective classes}, Comment. Math. Helv. 100 (2025), 269--322,
\href{https://ems.press/content/serial-article-files/50355}{published text};
\href{https://arxiv.org/pdf/2302.02541v3}{arXiv v3} distinguished.
\bibitem{DDNLmetric} T. Darvas, E. Di Nezza and C. H. Lu,
\emph{The metric geometry of singularity types},
\href{https://arxiv.org/pdf/1909.00839}{primary preprint},
Lemma 3.4 and Theorem 3.3, in its stated metric normalization.
\bibitem{DZ} T. Darvas and K. Zhang, \emph{Twisted K\"ahler--Einstein
metrics in big classes},
\href{https://arxiv.org/pdf/2208.08324v4}{arXiv:2208.08324v4},
Theorems 3.7 and 3.9 and their analytic proof units.
\bibitem{BT} E. Bedford and B. A. Taylor, \emph{A new capacity
for plurisubharmonic functions}, Acta Math. 149 (1982), 1--40,
\href{https://www.math.stonybrook.edu/~ebedford/PapersForM655/BTNewCapacity.pdf}
{primary text}, Theorem 2.1 and Corollary 2.2.
\bibitem{DDNL} T. Darvas, E. Di Nezza and C. H. Lu,
\emph{Monotonicity of non-pluripolar products and complex
Monge--Amp\`ere equations with prescribed singularity},
\href{https://arxiv.org/pdf/1705.05796v3}{arXiv:1705.05796v3},
Theorems 2.4 and 3.12.
\bibitem{PSWZ} K. Pang, H. Sun, Z. Wang and X. Zhou,
\emph{Capacity stability of complex Monge--Amp\`ere equations with
moving prescribed singularities},
\href{https://arxiv.org/pdf/2607.12797v1}{arXiv:2607.12797v1},
14 July 2026; comparison only, not an applied premise here.
\bibitem{Stacks} The Stacks Project Authors,
\href{https://stacks.math.columbia.edu/tag/0D2S}{Lemma 37.50.3,
tag 0D2S}, fiberwise criterion for relative ampleness for a proper
morphism; its proof is checked in the primary text.
\end{thebibliography}
\end{document}
